\section{Confinement}
\label{sec:sandbox}

Agent code is the third adversary (A3): written carelessly, written to
harm, or steered by what it read. The sandbox keeps five promises, each
by a separate mechanism so that a gap in one is not a gap in all: an
agent (P1)~reads and writes only its own files, (P2)~runs as its own
user, (P3)~reaches only the hosts its manifest named, (P4)~uses bounded
resources, and (P5)~cannot make the system calls escapes are built
from.

\subsection{Topology}

Agent code runs only in the \emph{agents' container}: the runtime's
image started with another command (\file{ai\_runtime/agents\_container/}).
Its one platform process, the \emph{spawner}, accepts structured,
key-authenticated requests from the runtime---check confinement,
start a worker, build requirements, read files, stop, report usage---and
relays worker protocol lines unread. The
container's network is internal and reaches only the runtime: no route
to the internet, the database, or the backend. It is handed no platform
secret. There is no unconfined path: a runtime without a spawner, or
whose spawner cannot confine, raises \code{ConfinementFailed} and does
not start. Only the test suite starts agent code without the agents'
container, through a double that exists in the tests alone.

\subsection{The spawn helper}

Workers are started by a small C program
(\file{ai\_runtime/agents\_container/spawn\_helper.c}), executable only
by the runtime's group. It is \emph{not} setuid-root: the file carries
exactly four capabilities (\code{cap\_setuid}, \code{cap\_setgid},
\code{cap\_chown}, \code{cap\_fowner}). It is executed in place of the
worker, on the same pipes, and applies, in order (\cref{fig:layers}):

\begin{figure}[t]
\centering
\begin{tikzpicture}[font=\small,
  st/.style={draw=accent,rounded corners=2pt,fill=white,minimum width=0.86\columnwidth,
             minimum height=1.55em,align=left,inner sep=3pt,text width=0.82\columnwidth},
  arr/.style={-{Stealth[length=4pt]},draw=ink!70}]
  \node[st,fill=soft] (a) {\textbf{0} refuse wrong caller, uid outside the worker range, home not owned by uid};
  \node[st,below=0.32em of a] (b) {\textbf{1} \code{setrlimit}: processes, open files, file size, no core};
  \node[st,below=0.32em of b] (c) {\textbf{2} \code{prctl(PR\_SET\_NO\_NEW\_PRIVS)}};
  \node[st,below=0.32em of c] (d) {\textbf{3} \code{setgroups}, \code{setresgid}, \code{setresuid}; \code{capset} $\emptyset$};
  \node[st,below=0.32em of d] (e) {\textbf{4} Landlock ruleset at the kernel's ABI};
  \node[st,below=0.32em of e] (f) {\textbf{5} seccomp filter (\code{prctl(PR\_SET\_SECCOMP, SECCOMP\_MODE\_FILTER)})};
  \node[st,below=0.32em of f,fill=soft] (g) {\textbf{6} \code{execve} the worker};
  \foreach \x/\y in {a/b,b/c,c/d,d/e,e/f,f/g} \draw[arr] (\x) -- (\y);
\end{tikzpicture}
\caption{The spawn helper's sequence. Each later step can only remove
access; nothing after step~3 can regain a privilege.}
\label{fig:layers}
\end{figure}

\code{no\_new\_privs} is what makes a capability-bearing binary safe to
leave on the worker's path: under it, an \code{exec} grants no file
capabilities, so a worker that runs the helper gains nothing. After the
uid switch the capability set is emptied, so no capability survives
into the worker. Landlock is applied after the switch and seccomp last,
so the helper's remaining work needs nothing the filter refuses. CPU
time is deliberately not limited by \code{RLIMIT\_CPU}, which counts a
process's whole life while workers are kept warm; each call is bounded
by the runtime's clock instead. Further subcommands chown, empty, and
sweep a uid's places, and \code{stop} kills every process of a uid, so
a browser never outlives the worker that started it.

\paragraph{One uid per approval}
Each approved agent gets a uid from a reserved range (one is reserved
for verifying packages, one for the builder). A worker's home is
emptied before every start, so it holds no state between starts.

\paragraph{Why not namespaces}
Inside a container's default profile a process cannot create user
namespaces; permitting it would mean relaxing the container's own
seccomp profile for everything in it. Separate uids, Landlock and a
seccomp filter provide the isolation without that cost.

\subsection{Files: Landlock}

The spawner names the paths; the helper hands them to the
kernel~\cite{landlock}. Read-and-execute: system directories, the
image's browsers, the agent's own package and environment. Read-write:
its home and spool. Write where ownership still decides: \code{/dev},
\code{/proc}, \code{/tmp}. Nothing else---not the platform's code,
other agents' packages, or the uid table. The helper validates none of
the paths: a Landlock ruleset can only remove access. The ruleset
requests what the running kernel supports (\cref{tab:landlock}). ABI~6
closes a channel that ownership and the firewall both miss: an abstract
Unix socket has a name but no inode, owner, or mode, and is invisible
to an address-based rule; from ABI~6 a fenced worker connects only to
abstract sockets created inside its own Landlock domain.

\begin{table}[t]
\centering\small
\caption{Landlock features requested by kernel ABI.}
\label{tab:landlock}
\begin{tabular}{@{}cll@{}}
\toprule
\textbf{ABI} & \textbf{Linux} & \textbf{Adds} \\
\midrule
1 & 5.13 & path-based file access \\
2 & 5.19 & \code{REFER}: moving files between own folders \\
3 & 6.2 & \code{TRUNCATE} \\
4 & 6.7 & TCP connect only to the proxy's port \\
6 & 6.12 & abstract Unix sockets scoped to own domain \\
\bottomrule
\end{tabular}
\end{table}

\subsection{System calls: seccomp}

The filter~\cite{seccomp} is a \emph{denylist}, on purpose: an agent may
bring any Python package, and an allowlist would be a list of what
somebody forgot. \Cref{tab:seccomp} shows what is refused for the
worker and all its descendants. \code{clone3} returns \code{ENOSYS}
rather than \code{EPERM}: its flags sit behind a pointer a BPF filter
cannot dereference, and glibc, told the call does not exist, falls back
to \code{clone}, whose flags the filter can inspect. Refusals return
the error codes programs already handle, so a package that probes a
call degrades instead of crashing. The filter is independent of the
container engine's profile and refuses what Docker's default allows,
\code{ptrace} among it. It is written so that the browser agents rely
on keeps working under it.

\begin{table}[t]
\centering\small
\caption{System calls refused to every worker.}
\label{tab:seccomp}
\begin{tabularx}{\columnwidth}{@{}>{\raggedright\arraybackslash}p{0.24\columnwidth}>{\raggedright\arraybackslash}X@{}}
\toprule
\textbf{Class} & \textbf{Calls} \\
\midrule
namespaces, mounts & \code{unshare}, \code{clone} with namespace flags, \code{setns}, \code{mount}, \code{umount2}, \code{pivot\_root}, the new mount API \\
cross-process memory & \code{ptrace}, \code{process\_vm\_readv/writev}, \code{pidfd\_getfd} \\
large kernel surface & \code{bpf}, \code{perf\_event\_open}, \code{userfaultfd}, \code{io\_uring\_*}, module and kexec loading, keyrings, \code{open\_by\_handle\_at} \\
socket families & all but \code{AF\_UNIX}, \code{AF\_INET}, \code{AF\_INET6}, and netlink for routes and device events \\
\bottomrule
\end{tabularx}
\end{table}

\subsection{Network: proxy, rule, hosts}

The egress proxy (\file{ai\_runtime/agents/egress.py}), running in the
runtime, (1)~identifies the worker by a per-worker pass and serves no
one else; (2)~refuses a host the manifest does not name; (3)~resolves
the name itself and refuses the host if \emph{any} answer is internal
(loopback, private, link-local---where cloud metadata lives), which
closes DNS rebinding toward the platform; and (4)~connects to the
address it checked and copies bytes. It terminates no TLS and installs
no certificate. Connections are bounded per worker and in total, and
every connection made or refused is recorded with its owner, host,
bytes each way, and duration. A host named by a credential or allowed
on a code card is \emph{lent} to the worker's pass only for the call.

The proxy is enforceable only because nothing else is reachable. At
container start, a root step installs an nftables rule on the output
hook matching the \emph{socket owner's uid} across the worker range:
accept TCP to one loopback port---the spawner's, which passes what
arrives to the runtime's proxy---and reject everything
else---other ports, other agents, the spawner, DNS. It then drops to
the platform's user with \code{net\_admin} removed from the bounding
set, so nothing started later can change the rule. Landlock's TCP rule
(ABI~4) is a second hold behind it.

\subsection{Dependencies: the builder}

A source distribution's build is arbitrary code. Installation is
therefore split: a dedicated \emph{builder} uid, under the same limits
and fences, whose proxy pass opens only the package index, downloads
and builds the declared requirements into wheels; one build at a time,
killed after five minutes, nothing persisting between builds. The
runtime then unpacks the wheels: it is handed files, never the list,
resolves and downloads nothing, and unpacking a wheel runs none of its
code.

\subsection{Memory}

A per-process address-space limit breaks browsers, so agents share a
memory budget. The spawner samples each agent's unique memory every
second and, at 90\% of the container's limit, ends the largest agent
with everything it runs, telling the affected call why.

\subsection{The inspector}

A sandbox one cannot see must be taken on faith. At every start,
before serving anyone, the platform tests its own locks: it starts a
process as an agent's user, asks the kernel which fences it can
enforce, prepares a place end to end, and tries to reach the network as
a worker would. It reports in plain words---\emph{workers are confined},
\emph{files are fenced}, \emph{system calls are filtered}---or which is
not, and why. If agents cannot be confined at all, the platform does not
start. If one lock is weaker on this machine (an older kernel without
Landlock, say), it runs, and every agent's page says what is not
enforced where it runs.
